\section{Conclusion}
\label{sec:conclusion}

This whitepaper has presented the empirical foundation of the PolieBotics Truth Beam, the projector--camera authentication system whose apparatus and methods are described in patent~\cite{patent2025}. The central claim the work establishes is narrow but, within its scope, strongly supported: a genuine recording made under the protocol carries a \emph{learnable physical chain-coupling} between each raw capture $C_t$ and the chain-derived emission $E_t$, and this coupling is sufficient to separate authentic captures from substrate-altered ones at near-perfect accuracy on the evaluated corpus. We close by recapitulating what the evidence does and does not support, and by situating the results relative to the patented system.

\subsection{What is established}

\paragraph{A near-perfect chain-coupling verifier.}
The protocol expands a \SI{32}{byte} chain state $S_t$ via BLAKE3-XOF~\cite{blake3} into a deterministic fBm emission tile $E_t$ (\num{43110} bytes per channel, four octaves, bit-exact integer rendering), projects it, and folds the BLAKE3 hash of the raw \SI{8}{bit} BayerRG8 capture $C_t$ back into $S_{t+1}$, closing a tamper-evident feedback loop in which any substituted capture induces a divergent downstream emission. The Phase~G verifier --- a custom \num{39.77}M-parameter $\epsilon$-prediction U-Net with ControlNet-style hint injection~\cite{ho2020ddpm,nicholdhariwal2021,zhang2023controlnet} --- learns this coupling from real same-rig pairs and scores a capture by its diffusion residual $\lVert \hat{\epsilon}-\epsilon\rVert^2$ at five fixed evaluation timesteps, never denoising. On held-out blocks of both recorded sessions it attains AUROC~$=1.000$ separating correct-$E$ from offset/wrong-$E$ captures (\Cref{fig:roc}), while an architecture- and data-matched shuffled control collapses to chance (AUROC~$\approx 0.5006$, $\Delta_{\text{wrong}}\approx 10^{-7}$), ruling out generic emission statistics and architectural artifact as the source of the signal. At the visualization tier the same substrate yields frame-level AUROC~$=0.9999$ (correct vs.\ shuffled $E$) and $0.9998$ (correct vs.\ synthetic $C$ at attacker step \num{100000}) across the \num{9735}-frame corpus (\Cref{fig:main_results}). The signal is real, reproducible, and not an inadvertent consequence of the emission's marginal distribution.

\paragraph{Same-rig cross-session and cross-protocol generalization.}
The verifier generalizes across recording sessions of the same rig and across two distinct chain protocols: D2 was recorded under the v9 chain and V10 under the v10 chain (the latter additionally committing a \texttt{ai\_payload\_root} for AI-agent response frames). A verifier trained on D2 alone reaches AUROC~$=1.000$ on held-out V10, and a V10-only verifier reaches AUROC~$=1.000$ on held-out D2, for the real-correct vs.\ shuffled/wrong-$E$ task. We emphasize that this is \emph{same-rig} cross-session (and cross-protocol) generalization only; it does not establish, and we do not claim, cross-rig, cross-camera, cross-projector, or cross-subject generalization. The harder fake-pair discrimination degrades out-of-distribution (e.g.\ a D2-only verifier drops to AUROC~$=0.8588$ on held-out V10 fakes), confirming that part of the learned coupling is session-specific.

\paragraph{F-A v1 is defeated.}
Against the learned adversary F-A v1 --- a \num{42.3}M-parameter ControlNet-style editor that synthesizes fake captures conditioned on a target emission --- Phase~G achieves real-vs-fake AUROC~$=1.000$ at every evaluated checkpoint (\num{5000}, \num{25000}, \num{70000}, \num{100000}) in both sessions, and the attacker signature does not converge as training proceeds (\Cref{fig:attacker_progression}). Crucially, this is not a missing-target-$E$ artifact: by step \num{100000}, F-A v1 has learned measurable target-$E$ coupling (fake-correct vs.\ fake-shuffled AUROC $\approx 0.99$ on D2 and $\approx 0.93$ on V10), so it is caught for failing the real-capture conditional distribution, raw-capture realism, and source-leakage --- not for lacking emission structure. Logistic probes on the diffusion features saturate at AUROC~$=1.000$, including under leave-one-checkpoint-out cross-validation, confirming that F-A v1 is, by these measures, an easy target.

\paragraph{The discriminative coupling is localized off-body.}
Under proper Mask R-CNN~\cite{he2017maskrcnn} body segmentation, preserving only off-body content (E3$_{\mathrm{seg}}$) yields AUROC~$=1.000$ across all twelve session$\times$condition cells, while preserving only the body (E2$_{\mathrm{seg}}$) drops to chance or below (\numrange{0.009}{0.538}; see \Cref{fig:segmentation}). The off-body region --- background scene, projection on static surfaces, sensor frame edges --- carries the chain-coupled signal that F-A v1 fails to reproduce, because the attacker smooths body content and does not perturb off-body content in a chain-coupled way. The body-region signal is not absent: it survives at low spatial frequencies (a Gaussian-blurred residual at $\sigma=4$ retains AUROC~$\approx 0.997$--$0.998$ on E2$_{\mathrm{seg}}$) that the high-frequency-dominated $\epsilon$-MSE washes out via the likelihood-overspread failure mode of generative out-of-distribution detectors~\cite{nalisnick2019}. We therefore frame the body-only result in distribution-distance terms rather than as ``fakes look more real than real.'' The orthogonal per-frame relative-comparison test (EXP~6: top-1~$=100\%$ across 120 frames, 51 candidates each) remains the load-bearing chain-coupling demonstration and is unaffected by this absolute-score inversion.

\paragraph{Scoring-function and sensitivity structure.}
The verifier's sensitivity to emission perturbations is graded and octave/persistence dependent. Gross perturbations --- full replacement, channel swaps, low-octave swaps, large frame offsets --- are perfectly separable (AUROC~$=1.000$). Fine XOF bit-flips (Type~1 small-$k$ global flips, Type~3 localized flips) sit near chance (AUROC~$\approx 0.50$), with separability emerging only at large $k$ and in low-frequency octaves; these responses must be read against the rendered-emission difference $\Delta E$, since very small bit-flips fall below the optical-relevance threshold of the \SI{8}{bit}, optically low-pass channel. The supervised Phase~H baseline is E-aware for coarse contrasts (target vs.\ shuffled and target vs.\ source both AUROC~$=1.000$) but, owing to a per-frame label-locking limitation in its training regime, is blind to fine XOF perturbations; it is a coarse supervised baseline, not a production verifier.

\subsection{What remains open}

\paragraph{The adaptive attacker.}
The strongest claim a forensic verifier can make --- robustness against an adversary that explicitly optimizes against it --- is \emph{not} established here. The stronger attacker F-A v2, specified to be trained against fresh binder surrogates and to perturb off-body regions in a chain-coupled manner, exists only as a design specification; its trainer code does not exist and it has not been trained. Phase~G is deliberately held out as a final-evaluation oracle and must never enter F-A v2 training, validation, tuning, or selection. Whether the off-body coupling signature survives an attacker that targets it directly is the central open question, and adaptive-attacker robustness is strictly future work.

\paragraph{Cross-rig generalization.}
All results rest on a same-rig, single-operator, two-session corpus (D2 and V10) featuring one human performer. No cross-rig, cross-camera, cross-projector, or cross-subject experiment exists. Whether the learned coupling --- or a verifier trained to recognize it --- transfers to a different optical assembly or a different scene is untested, and the partial session-specificity already observed in the cross-session fake-pair degradation counsels against assuming such transfer without retraining or recalibration.

\paragraph{Calibration and protocol details.}
Several practical quantities remain unmeasured: the empirically optimal verifier-side frame-to-chain offset (the recording convention itself is offset $0$, $E_t = \mathrm{gen}(\mathrm{XOF}(S_t))$; what is unvalidated is the best alignment for the optical verifier on held-out data), the measured black level (recorded as $0$ with source ``no\_measurement\_found,'' with subtraction forbidden until confirmed), and the projector-pipeline lock status. These should be surfaced in any release manifest and resolved before any offset- or calibration-dependent operating point is deployed.

\subsection{Relation to the patent}

The patented apparatus~\cite{patent2025} specifies a projector--camera system that emits a structured illumination pattern and captures the scene under that pattern, with the projected pattern and the captured frame bound together so that tampering is detectable. The present work instantiates that apparatus concretely --- a BLAKE3-XOF~\cite{blake3} chain producing fBm emissions, a Sony IMX540 sensor delivering \SI{8}{bit} BayerRG8 captures at a \SI{2.5}{\hertz} chain-commit rate, and a drand-anchored~\cite{drandbeacon} tamper-evident state chain --- and supplies the missing empirical link: a demonstration that the physical coupling the patent relies upon is not merely cryptographically asserted but \emph{statistically recoverable} from real recordings, to the point of near-perfect verification within the evaluated scope. By construction of the cryptographic chain, a substituted capture yields a divergent downstream emission; the diffusion verifier demonstrates that the resulting inconsistency is detectable in the optical record itself, even against a learned forger that has access to the target emission. Together they convert the patent's tamper-evidence claim from a structural property of the protocol into a measured, defended property of the captured signal --- bounded, for now, to a single rig, two sessions, and a non-adaptive attacker, with the path to a stronger threat model laid out as future work.
